\chapter{Genetics}

\medterm{Allelic Heterogeneity}-- Phenomenon where different mutations at the same gene locus cause the same disease, often with variable expressivity or severity. For example, over 2,000 CFTR mutations cause cystic fibrosis with phenotypes ranging from mild (isolated congenital absence of the vas deferens) to severe (classic multisystem CF).

\medterm{Aneuploidy}-- An abnormal number of chromosomes in a cell, differing from the wild-type diploid number (46,XX or 46,XY). Caused primarily by meiotic non-disjunction during gametogenesis. Common autosomal trisomies: Down syndrome (trisomy 21), Edwards syndrome (trisomy 18), Patau syndrome (trisomy 13). Sex chromosome aneuploidies: Turner syndrome (45,X0), Klinefelter syndrome (47,XXY).

\medterm{Autosomal Dominant}-- A pattern of inheritance in which a single copy of a mutated gene on a non-sex chromosome (autosome) is sufficient to cause disease. Each child of an affected parent has a 50\% chance of inheriting the mutation. Characteristics: vertical transmission across generations, equal male and female frequency. Examples: Marfan syndrome, Huntington disease, Achondroplasia, Familial Hypercholesterolemia.

\medterm{Autosomal Recessive}-- A pattern of inheritance requiring two copies of a mutated gene (one from each heterozygous carrier parent) for disease manifestation. Carriers are clinically asymptomatic. Each child of two carrier parents has a 25\% chance of disease, 50\% chance of carrier status, and 25\% chance of non-carrier. Characteristics: horizontal transmission (affected siblings), increased frequency in consanguineous marriages. Examples: Cystic Fibrosis, Sickle Cell Anemia, Tay-Sachs Disease, Phenylketonuria.

\medterm{Chromosomal Abnormality}-- Any deviation from the normal number (46 chromosomes, 23 pairs) or structure of chromosomes. Numerical abnormalities: aneuploidy (trisomy 21/Down syndrome, trisomy 18/Edwards syndrome, trisomy 13/Patau syndrome, monosomy X/Turner syndrome), polyploidy (triploidy, tetraploidy). Structural abnormalities: deletions (Cri-du-chat, 22q11.2 deletion), duplications, translocations (Robertsonian, reciprocal), inversions, and ring chromosomes. Mosaicism: presence of two or more genetically different cell lines in the same individual. Screening: combined test (nuchal translucency, hCG, PAPP-A), NIPT (cell-free fetal DNA), quadruple test. Diagnosis: chorionic villus sampling (11-14 weeks), amniocentesis (15+ weeks). Karyotype, FISH, or chromosomal microarray confirms the diagnosis.

\medterm{Chromosomes}-- The thread-like structures within the cell nucleus that carry genetic information in the form of DNA. Humans have 46 chromosomes arranged in 23 pairs: 22 pairs of autosomes and 1 pair of sex chromosomes (XX in females, XY in males). Each parent contributes one chromosome to each pair. Chromosomes are visible under a light microscope during cell division and are arranged in a standardized format called a karyotype (ordered by size, centromere position, and banding pattern, using G-banding or spectral karyotyping). The short arm is designated p (petit) and the long arm is q. Chromosomal abnormalities can involve number (aneuploidy) or structure (deletions, duplications, translocations, inversions). Telomeres: protective caps at chromosome ends. Centromere: region where sister chromatids are joined and spindle fibers attach during cell division.

\medterm{CRISPR-Cas9}-- A bacterial adaptive immune mechanism adapted as a targeted genome editing tool. Uses a synthetic single guide RNA (sgRNA) to direct the Cas9 endonuclease to a specific target DNA sequence adjacent to a protospacer adjacent motif (PAM). Introduces double-strand breaks repaired by non-homologous end joining (NHEJ, gene knockout) or homology-directed repair (HDR, precise insertion/edit).

\medterm{Cystic Fibrosis}-- An autosomal recessive multisystem disorder caused by mutations in the CFTR gene on chromosome 7q31.2, encoding the cystic fibrosis transmembrane conductance regulator, a cAMP-regulated chloride channel expressed on epithelial surfaces. Over 2,000 CFTR mutations have been identified; the most common is F508del (deletion of phenylalanine at position 508), accounting for approximately 70 percent of alleles worldwide. CFTR mutations are classified by defect type: Class I (no protein production -- nonsense, frameshift); Class II (defective trafficking and degradation -- F508del); Class III (defective gating -- G551D); Class IV (defective chloride conductance); Class V (reduced synthesis); Class VI (accelerated turnover). Clinical features: chronic suppurative lung disease (bronchiectasis, Pseudomonas aeruginosa, progressive obstructive airways disease), exocrine pancreatic insufficiency (steatorrhea, fat-soluble vitamin deficiency), meconium ileus in neonates, congenital bilateral absence of the vas deferens (male infertility), and elevated sweat chloride (>60 mmol/L, measured by pilocarpine iontophoresis). Management: chest physiotherapy, mucolytics (dornase alfa), inhaled antibiotics (tobramycin), pancreatic enzyme replacement, fat-soluble vitamin supplementation, and CFTR modulator therapy (ivacaftor for Class III; lumacaftor/ivacaftor and tezacaftor/ivacaftor for Class II; elexacaftor/tezacaftor/ivacaftor triple therapy). Lung transplantation in end-stage disease.

\medterm{Epigenetics}-- Heritable changes in gene expression or cellular phenotype caused by mechanisms other than alterations in the underlying DNA sequence. Mechanisms include DNA methylation (typically silences gene expression at CpG islands), histone modification (acetylation relaxes chromatin to activate transcription; methylation can activate or repress), and non-coding RNA regulation. Influenced by environmental factors, aging, and diet.

\medterm{Friedreich Ataxia}-- The most common inherited ataxia, an autosomal recessive disorder caused by GAA trinucleotide repeat expansion in the FXN gene on chromosome 9q13, encoding frataxin, a mitochondrial protein essential for iron-sulfur cluster biogenesis. The repeat expansion (typically 66-1,700 repeats; normal is 5-30) causes gene silencing through heterochromatin formation, leading to frataxin deficiency (20-30 percent of normal levels). Frataxin deficiency impairs mitochondrial iron metabolism, causing iron accumulation within mitochondria, defective iron-sulfur cluster assembly, impaired oxidative phosphorylation, and increased oxidative stress. Clinical triad: progressive ataxia (cerebellar and sensory), hypertrophic cardiomyopathy (most common cause of premature death), and diabetes mellitus or glucose intolerance. Additional features: dysarthria, lower limb areflexia, loss of vibration and proprioception (dorsal columns), pyramidal weakness, pes cavus, scoliosis, and optic atrophy. Age of onset typically 5-25 years; earlier onset correlates with larger repeat expansions. Diagnosis: genetic testing for GAA expansion in the FXN gene. Management is supportive: physical therapy, speech therapy, cardiac monitoring (annual echocardiogram, ECG), diabetes screening. No disease-modifying therapy available; omaveloxolone approved in some populations. Prognosis: loss of ambulation 10-20 years after onset; death usually from cardiomyopathy.

\medterm{Gamete}-- A mature haploid reproductive cell (23 chromosomes) that fuses with another gamete during fertilization to form a diploid zygote. Male gamete: spermatozoon (flagellated, motile, produced by spermatogenesis in testes, approximately 100 million per day). Female gamete: oocyte (ovum, non-motile, larger, produced by oogenesis in ovaries; approximately 1 million present at birth, reduced to 400-500 ovulated in a lifetime). Gametogenesis involves meiosis (reduction division) to achieve haploidy. Abnormal gametes may occur from meiotic non-disjunction, leading to aneuploidy in offspring (Down syndrome, Turner syndrome, Klinefelter syndrome).

\medterm{Gaucher Disease}-- Most common lysosomal storage disorder caused by glucocerebrosidase deficiency. Results in accumulation of glucocerebroside in macrophages. Type 1 is the most common and affects the liver, spleen, and bones.

\medterm{Gene}-- The fundamental physical and functional unit of heredity. A gene is a segment of DNA (or RNA in some viruses) that encodes a functional product (typically a protein or, in the case of non-coding RNA genes, a functional RNA molecule). The human genome contains 20,000-25,000 protein-coding genes. Structure: promoter region (regulatory sequence upstream, binds RNA polymerase and transcription factors), exons (coding sequences), introns (non-coding intervening sequences, spliced out), and untranslated regions (UTRs, 5 and 3 prime). Gene expression is regulated at multiple levels: transcriptional (transcription factors, enhancers, silencers, epigenetic modifications), post-transcriptional (splicing, RNA editing, mRNA stability), translational (initiation factors, miRNA, ribosome binding), and post-translational (phosphorylation, ubiquitination, glycosylation). Gene nomenclature: official symbols assigned by HGNC (HUGO Gene Nomenclature Committee), e.g., CFTR, BRCA1, TP53.

\medterm{Genetic Code}-- The set of rules by which information encoded in nucleotide sequences (DNA and RNA) is translated into proteins. The code is: (1) Triplet: three consecutive nucleotides (codon) specify one amino acid. (2) Degenerate: 64 possible codons encode 20 amino acids and 3 stop codons (UAA, UAG, UGA). (3) Universal: nearly identical across all organisms. (4) Non-overlapping: each nucleotide belongs to only one codon. (5) Commaless: codons are read sequentially without gaps. Key features: start codon (AUG, methionine, also sets reading frame), wobble hypothesis (third base pairing is flexible, allowing fewer tRNAs to decode all codons).

\medterm{Genetic Engineering}-- The direct manipulation of an organism's genes using biotechnology. Techniques include: (1) Recombinant DNA technology: insertion of foreign DNA into a vector (plasmid, virus) for expression in a host organism (bacteria, yeast, mammalian cells). (2) CRISPR-Cas9: RNA-guided nuclease that cuts DNA at a specific sequence, enabling gene knockout, knock-in, and precise editing. (3) Gene silencing: RNA interference (siRNA, shRNA) to reduce gene expression. (4) Gene therapy: delivery of functional genes to treat genetic disorders. Applications: production of therapeutic proteins (insulin, growth hormone, clotting factors, monoclonal antibodies), genetically modified organisms (GMOs, crop improvement), gene therapy (SCID, hemophilia, sickle cell disease, Leber congenital amaurosis), synthetic biology, and xenotransplantation. Ethical and safety concerns: off-target effects, germline editing consequences, ecological impact of GMOs.

\medterm{Genetics}-- The branch of biology concerned with the study of genes, genetic variation, and heredity in living organisms. Major subdisciplines: (1) Mendelian genetics: patterns of inheritance of single-gene traits (autosomal dominant, autosomal recessive, X-linked, Y-linked, mitochondrial). (2) Molecular genetics: structure and function of DNA, RNA, and proteins at the molecular level. (3) Population genetics: allele frequency distribution and change under the influence of natural selection, genetic drift, mutation, and gene flow (Hardy-Weinberg equilibrium). (4) Quantitative genetics: inheritance of complex traits influenced by multiple genes and environment (heritability estimates). (5) Genomics: the comprehensive study of entire genomes (structural, functional, comparative, epigenomics). (6) Clinical genetics: diagnosis and management of genetic disorders, genetic counseling, prenatal diagnosis, and carrier screening.

\medterm{Genome}-- The complete set of genetic material (DNA) of an organism. Human genome: approximately 3.2 billion base pairs, 20,000-25,000 protein-coding genes (only 1-2\% of total DNA), organized into 23 pairs of chromosomes (22 autosomes + sex chromosomes). Non-coding DNA (98\%) includes: regulatory sequences (promoters, enhancers), transposable elements, introns, telomeres, centromeres, pseudogenes, microRNAs, long non-coding RNAs, and repetitive sequences. The Human Genome Project (1990-2003) sequenced the human genome. Genomic medicine uses genome sequencing for diagnosis, pharmacogenomics, and prediction of disease risk.

\medterm{Genotype}-- The genetic constitution of an organism, determined by the specific alleles present at a given genetic locus. The genotype is the inherited genetic information carried by an individual, as distinct from the phenotype (observable physical or biochemical characteristics). Genotype is influenced by genotype-environment interactions. In Mendelian genetics, a homozygous genotype has two identical alleles at a locus (e.g., AA or aa), while a heterozygous genotype has two different alleles (Aa). Genotype frequency in a population is determined by Hardy-Weinberg equilibrium. Genotyping methods: PCR, DNA sequencing, microarray, SNP analysis.

\medterm{Hardy-Weinberg Equilibrium}-- A fundamental principle of population genetics stating that allele and genotype frequencies remain constant across generations in the absence of evolutionary influences (mutation, selection, migration, genetic drift, non-random mating). Equations: $p + q = 1$ (allele frequencies) and $p^2 + 2pq + q^2 = 1$ (genotype frequencies), where $2pq$ represents carrier frequency for autosomal recessive diseases.

\medterm{Hemochromatosis}-- Hereditary hemochromatosis (HH) is an autosomal recessive disorder of iron metabolism caused most commonly by homozygosity for the C282Y mutation in the HFE gene on chromosome 6p21.3, leading to inappropriately low hepcidin production and excessive intestinal iron absorption (1.5-2 mg/day normally, increasing to 3-4 mg/day in HH). Iron accumulates progressively in parenchymal organs (liver, pancreas, heart, pituitary, joints, skin), causing tissue damage through oxidative stress (Fenton reaction generating hydroxyl radicals that damage lipids, proteins, and DNA). Clinical features: bronze or gray skin pigmentation (bronze diabetes), diabetes mellitus, cirrhosis (increased risk of hepatocellular carcinoma), dilated or restrictive cardiomyopathy, hypogonadotropic hypogonadism, arthritis (typically second and third metacarpophalangeal joints), and osteoporosis. Penetrance is incomplete; many homozygotes show only biochemical abnormalities without clinical disease. Diagnosis: elevated transferrin saturation (>45 percent) and serum ferritin; HFE genetic testing confirms; liver MRI or biopsy quantifies iron burden. Treatment: graded phlebotomy (500 mL weekly until ferritin <50 mcg/L, then maintenance every 2-4 months). Iron chelation (deferasirox, deferoxamine) for patients who cannot tolerate phlebotomy. First-degree relatives should be screened with iron studies and HFE genotyping.

\medterm{Hemophilia A}-- X-linked recessive bleeding disorder caused by deficiency of coagulation factor VIII. Severity correlates with factor activity level. Characterized by hemarthrosis, soft tissue bleeding, and prolonged aPTT.

\medterm{Hemophilia B}-- X-linked recessive bleeding disorder caused by deficiency of coagulation factor IX. Also known as Christmas disease. Clinically indistinguishable from hemophilia A. Diagnosed by factor IX activity assay.

\medterm{Huntington Disease}-- An autosomal dominant neurodegenerative disorder caused by an expanded CAG trinucleotide repeat in the HTT gene on chromosome 4p16.3, encoding huntingtin protein. Normal alleles have fewer than 26 CAG repeats; 27-35 repeats are intermediate (not pathogenic but may expand in offspring); 36-39 repeats show reduced penetrance; 40 or more repeats cause full penetrance disease. The expanded polyglutamine tract causes mutant huntingtin to misfold, aggregate, and disrupt cellular processes (transcription, mitochondrial function, axonal transport, synaptic transmission), leading to neuronal death (most prominently in the caudate nucleus, putamen, and cerebral cortex). Clinical triad: chorea (involuntary, dance-like movements; the most characteristic motor feature), progressive dementia (executive dysfunction, memory loss, visuospatial deficits), and psychiatric disturbance (depression (most common), anxiety, irritability, obsessive-compulsive symptoms, psychosis). Mean age of onset 35-44 years; juvenile-onset HD (CAG >60 repeats) presents with rigidity, bradykinesia, and seizures rather than chorea. Anticipation occurs -- the CAG repeat expands during male meiosis (paternal transmission), resulting in earlier onset in successive generations. Diagnosis: genetic testing for CAG repeat expansion in the HTT gene (with pre-test and post-test genetic counseling). Treatment is symptomatic: tetrabenazine or deutetrabenazine for chorea; antidepressants; antipsychotics. No disease-modifying therapy exists. Prognosis: mean survival 15-20 years from onset.

\medterm{Karyotype}-- The complete visual profile of an individual\'s chromosomes in metaphase, arranged in homologous pairs by size, centromere position, and banding pattern (Giemsa staining). Used to detect numerical chromosomal abnormalities (trisomies, monosomies) and large structural rearrangements (translocations, inversions, deletions).

\medterm{Li-Fraumeni Syndrome}-- Autosomal dominant cancer predisposition syndrome most commonly caused by TP53 mutations. Characterized by early-onset cancers including sarcomas, breast cancer, brain tumors, and adrenocortical carcinoma. Very high lifetime cancer risk.

\medterm{Marfan Syndrome}-- An autosomal dominant connective tissue disorder caused by mutations in the FBN1 gene on chromosome 15q21.1, encoding fibrillin-1, a major component of microfibrils that form the scaffold for elastin in connective tissue. Fibrillin-1 mutations also lead to increased transforming growth factor-beta (TGF-beta) signaling, contributing to the pathogenesis. The Ghent nosology (2010 revision) diagnoses Marfan syndrome when aortic root dilation (Z-score >=2) AND ectopia lentis are present; or aortic root dilation AND a pathogenic FBN1 mutation; or aortic root dilation AND a sufficient systemic score (points for wrist and thumb signs, pectus carinatum, hindfoot deformity, pneumothorax, dural ectasia, protrusio acetabuli, facial features, skin striae, myopia, mitral valve prolapse). Skeletal features: tall stature, arachnodactyly, pectus deformities, scoliosis, joint hypermobility. Cardiovascular: aortic root aneurysm (leading cause of mortality), aortic regurgitation, mitral valve prolapse, aortic dissection (type A). Ocular: ectopia lentis (upward lens dislocation), myopia, retinal detachment. Diagnosis: revised Ghent criteria; FBN1 genetic testing; serial echocardiography for aortic root measurement; slit-lamp ophthalmology. Management: beta-blockers (atenolol) and losartan to slow aortic root growth; serial echocardiography; annual ophthalmology review; avoidance of contact sports and heavy exertion. Elective prophylactic aortic root replacement (Bentall procedure) when aortic root diameter reaches 50 mm (or 45 mm with rapid progression or family history of dissection).

\medterm{Meiosis}-- A specialized form of cell division in germ cells that reduces the chromosome number by half, producing four genetically distinct haploid gametes (sperm or ova) from one diploid progenitor. Consists of two rounds: Meiosis I (reductional division; homologous chromosomes pair, undergo crossing over/recombination, and segregate) and Meiosis II (equational division; sister chromatids segregate).

\medterm{Mendelian Inheritance}-- Biological inheritance patterns obeying Gregor Mendel\'s laws: Law of Segregation (allele pairs separate during gamete formation) and Law of Independent Assortment (genes for different traits segregate independently). Includes autosomal dominant, autosomal recessive, and X-linked inheritance patterns.

\medterm{Mitosis}-- Process of somatic cell division resulting in two genetically identical diploid daughter cells from a single parent cell. Phases: Prophase (chromatin condenses), Prometaphase (nuclear envelope breaks down), Metaphase (chromosomes align at equatorial plate), Anaphase (sister chromatids separate to opposite poles), Telophase (nuclear envelopes reform), followed by Cytokinesis.

\medterm{Non-disjunction}-- The failure of homologous chromosomes (Meiosis I) or sister chromatids (Meiosis II or Mitosis) to separate properly during cell division. Results in daughter cells with abnormal chromosome numbers (nullisomy, monosomy, disomy, trisomy). Advanced maternal age (>35 years) is the primary risk factor for meiotic non-disjunction.

\medterm{Penetrance}-- The proportion of individuals carrying a specific disease-causing gene mutation who actually express the associated clinical phenotype. Expressed as a percentage. Complete penetrance (100\%): all mutation carriers show disease (e.g., Huntington disease). Incomplete/Reduced penetrance (<100\%): some carriers remain asymptomatic (e.g., BRCA1/BRCA2 hereditary breast cancer).

\medterm{Phenotype}-- The composite observable physical, physiological, biochemical, and behavioral traits of an organism, resulting from the interaction of its genotype with environmental influences.

\medterm{Phenylketonuria}-- An autosomal recessive aminoacidopathy caused by deficiency of phenylalanine hydroxylase (PAH), the enzyme that converts phenylalanine to tyrosine. PAH deficiency leads to toxic accumulation of phenylalanine and its metabolites (phenylketones: phenylacetate, phenyllactate, phenylpyruvate) in blood and brain, causing irreversible intellectual disability if untreated. PKU is caused by mutations in the PAH gene on chromosome 12q23.2, with over 1,000 mutations identified. Classic PKU: blood phenylalanine >1,200 micromol/L (20 mg/dL) on normal diet; mild PKU: 600-1,200 micromol/L; mild hyperphenylalaninaemia: 120-600 micromol/L (no treatment required). Newborn screening by tandem mass spectrometry (Guthrie test) detects elevated phenylalanine levels in all developed countries; prompt initiation of dietary therapy prevents neurological damage. Treatment: lifelong low-phenylalanine diet (phenylalanine-free medical formula and restricted natural protein); tetrahydrobiopterin (BH4, sapropterin) response in some mutation types; pegvaliase (PEGylated recombinant phenylalanine ammonia lyase) for adults with poorly controlled PKU. Maternal PKU: women with PKU must maintain strict phenylalanine control before conception and throughout pregnancy to prevent fetal damage (microcephaly, intellectual disability, congenital heart disease, low birth weight).

\medterm{Sickle Cell Disease}-- A group of inherited hemoglobinopathies caused by a point mutation (glutamic acid to valine at position 6) in the beta-globin gene (HBB) on chromosome 11, producing hemoglobin S (HbS). Homozygous HbSS is the most common and severe form. Under conditions of deoxygenation, HbS polymerizes, causing red blood cells to assume a rigid sickle shape. Sickled cells are less deformable, leading to vaso-occlusion, microvascular ischemia, and hemolytic anemia. SCD manifests through acute and chronic complications: vaso-occlusive crises (bone pain, dactylitis, priapism), acute chest syndrome (fever, chest pain, pulmonary infiltrates, hypoxia), stroke (overt or silent cerebral infarction), bacterial infection risk (functional asplenia, encapsulated organisms -- Streptococcus pneumoniae, Haemophilus influenzae), aplastic crisis (parvovirus B19-induced temporary arrest of erythropoiesis), splenic sequestration, pulmonary hypertension, leg ulcers, avascular necrosis of femoral head, and chronic kidney disease. Diagnosis: hemoglobin electrophoresis or high-performance liquid chromatography (HPLC) showing HbS (typically >80 percent in HbSS). Newborn screening (isoelectric focusing) universal in high-prevalence regions. Treatment: hydroxyurea (increases fetal hemoglobin HbF, reducing sickling); regular blood transfusion for stroke prevention; penicillin prophylaxis and pneumococcal vaccination (functional asplenia); pain management; acute chest syndrome management (antibiotics, transfusion, bronchodilators). hematopoietic stem cell transplantation is curative; gene therapy (LentiGlobin, CRISPR-based approaches) is emerging. Prognosis: mean survival 40-50 years in high-income countries.

\medterm{Sporadic}-- Occurring in isolated, random, or scattered instances, without a predictable pattern or clear cause. In medical genetics, sporadic cases of a disease occur in individuals with no family history of the condition, resulting from de novo (new) mutations rather than inherited germline mutations. Sporadic cancers arise from acquired (somatic) mutations. Examples: most cases of breast, ovarian, and colorectal cancers are sporadic. Many genetic syndromes can occur sporadically through de novo mutations (e.g., neurofibromatosis type 1, CHARGE syndrome, Cornelia de Lange syndrome). Sporadic cases have a low recurrence risk in siblings unless gonadal mosaicism is present.

\medterm{Tay-Sachs Disease}-- Autosomal recessive lysosomal storage disorder caused by hexosaminidase A deficiency. Results in accumulation of GM2 ganglioside in neurons causing progressive neurodegeneration. Most common in Ashkenazi Jewish populations.

\medterm{Teratogenesis}-- The process by which congenital malformations (birth defects) are produced in a developing fetus by exposure to teratogens—agents that interfere with normal embryonic or fetal development. Teratogenesis is most sensitive during the period of organogenesis (weeks 3--8 of gestation), when major organs are forming. Principles: (1) Susceptibility depends on the genotype of the embryo and the timing, dose, and duration of exposure. (2) Different teratogens produce characteristic patterns of malformations. (3) The critical period for each organ system is when it is undergoing rapid differentiation. Major teratogens: (1) Medications—thalidomide (limb defects—phocomelia, amelia), isotretinoin (CNS, cardiac, craniofacial defects), valproic acid (neural tube defects, learning difficulties), warfarin (nasal hypoplasia, stippled epiphyses), methotrexate (CNS, skeletal), ACE inhibitors (renal dysplasia, oligohydramnios), tetracycline (bone and tooth discoloration, bone growth suppression), and androgens/lithium/carbamazepine/phenytoin. (2) Maternal infections—rubella (deafness, cataracts, congenital heart disease), cytomegalovirus, toxoplasmosis, varicella, Zika virus. (3) Maternal conditions—diabetes (neural tube defects, congenital heart disease, macrosomia), phenylketonuria (microcephaly, intellectual disability). (4) Environmental exposures—alcohol (fetal alcohol syndrome—facial dysmorphism, growth restriction, neurodevelopmental impairment), ionizing radiation (microcephaly, intellectual disability), methylmercury (cerebral palsy, blindness), tobacco smoke (IUGR, cleft palate, preterm birth). Prevention: preconception counseling, folic acid supplementation (4--5 mg/day in high-risk women—reduces neural tube defects by 70\%), and avoidance of known teratogens during pregnancy.

\medterm{Thalassemia}-- Group of autosomal recessive disorders characterized by reduced or absent globin chain synthesis. Alpha thalassemia results from mutations in the HBA genes. Beta thalassemia results from mutations in the HBB gene.

\medterm{Trisomy}-- The presence of an extra chromosome (three copies instead of the normal pair), resulting in a total of 47 chromosomes. Trisomy is the most common type of aneuploidy and usually arises from non-disjunction during meiosis (failure of homologous chromosomes or sister chromatids to separate). Trisomies involving autosomes (non-sex chromosomes) are generally incompatible with life, except for trisomy 21, 18, and 13. Trisomy 21 (Down syndrome): most common viable autosomal trisomy (incidence 1 in 700 live births), associated with intellectual disability, congenital heart disease (AVSD, VSD, ASD), hypotonia, characteristic facial features (flat facies, upslanting palpebral fissures, epicanthal folds, protruding tongue), increased risk of leukemia (ALL), Alzheimer disease (by age 40), and atlantoaxial instability. Trisomy 18 (Edwards syndrome): incidence 1 in 5,000, severe intellectual disability, rocker-bottom feet, clenched hands with overlapping fingers, congenital heart disease, and omphalocele; median survival <1 year. Trisomy 13 (Patau syndrome): incidence 1 in 15,000, severe intellectual disability, cleft lip/palate, polydactyly, holoprosencephaly, and congenital heart disease; median survival 7--10 days. Sex chromosome trisomies: XXY (Klinefelter syndrome), XXX (triple X syndrome), XYY (Jacob syndrome)—generally milder phenotypes. Prenatal screening: combined test (nuchal translucency, $\beta$-hCG, PAPP-A) at 11--13 weeks, NIPT (cell-free fetal DNA), and diagnostic testing (CVS or amniocentesis).

\medterm{Wilson Disease}-- An autosomal recessive disorder of copper metabolism caused by mutations in the ATP7B gene on chromosome 13q14.3, encoding a copper-transporting ATPase. The defective ATP7B protein cannot incorporate copper into ceruloplasmin or excrete excess copper into bile, leading to toxic copper accumulation primarily in the liver and brain (basal ganglia, particularly the putamen and globus pallidus), with secondary deposition in the cornea (Kayser-Fleischer rings -- golden-brown copper deposits in Descemet's membrane at the limbus, visible on slit-lamp examination, pathognomonic for Wilson disease). Hepatic presentation: acute hepatitis, fulminant hepatic failure (with Coombs-negative hemolytic anemia), chronic hepatitis, cirrhosis, and portal hypertension. Neuropsychiatric presentation: movement disorders (tremor, dysarthria, dystonia, chorea, parkinsonism), dysphagia, personality changes, depression, psychosis, and cognitive decline. Other features: renal tubular dysfunction (Fanconi syndrome), arthropathy, and cardiomyopathy. Diagnosis: reduced serum ceruloplasmin (<0.2 g/L), elevated 24-hour urinary copper excretion (>1.6 micromol/24h), elevated hepatic copper on liver biopsy (>250 mcg/g dry weight); Kayser-Fleischer rings on slit-lamp examination; ATP7B genetic testing. The Leipzig diagnostic scoring system aids diagnosis. Treatment: lifelong copper chelation (D-penicillamine as first-line, trientine as second-line) and zinc acetate (blocks intestinal copper absorption, used as maintenance or for pre-symptomatic patients). Avoid copper-rich foods (shellfish, liver, nuts, chocolate). Liver transplantation for fulminant hepatic failure or decompensated cirrhosis. Prognosis: excellent with lifelong treatment; progressive and fatal without treatment.

\medterm{X-Linked Recessive}-- Inheritance of a gene mutation located on the X chromosome. Expressed in all hemizygous males ($XY$) and homozygous females ($XX$). Females carrying one mutant copy are asymptomatic carriers. Carrier mothers have a 50\% chance of passing disease to sons and 50\% chance of carrier status in daughters. Affected fathers pass the mutation to 100\% of daughters (carriers) and 0\% of sons. Examples: Hemophilia A/B, Duchenne Muscular Dystrophy, Glucose-6-Phosphate Dehydrogenase Deficiency.

